% =====================================================================
% Standalone report on the propagation factor F_star proposed in
% F_star/main_F_star.tex (the passages set in red).
%
% This file is self-contained. The preamble below copies the packages
% and the macro definitions from Paper/preamble.tex and Paper/main.tex
% that the report needs; it does not \input either file, and it does
% not read Paper/references.bib.
% =====================================================================
\documentclass[11pt]{article}

\usepackage[letterpaper, margin=1in]{geometry}
\usepackage{amsmath,amssymb}
\usepackage{booktabs}
\usepackage{xcolor}
\definecolor{darkgreen}{rgb}{0.0,0.5,0.0}
\usepackage[colorlinks,linkcolor=red,citecolor=darkgreen,urlcolor=blue]{hyperref}
\usepackage[amsmath,thmmarks]{ntheorem}

\newtheorem{remark}{Remark}
\newtheorem{lemma}{Lemma}
\newtheorem{theorem}{Theorem}
\newenvironment{proof}{\par\medskip\noindent\textit{Proof.}\ }{\hfill$\square$\par\medskip}
\theoremstyle{plain}

% Macros copied from Paper/preamble.tex
\newcommand{\Le}{\leqslant}
\newcommand{\Ge}{\geqslant}
\newcommand{\D}{\mathrm{d}}
\newcommand{\KL}{\mathrm{KL}}
\newcommand{\X}{\mathrm{X}}
\newcommand{\ESS}{\mathrm{ESS}}

% Abbreviations used in F_star/main_F_star.tex but not in Paper/preamble.tex
\newcommand{\osc}{\operatorname{osc}}
\newcommand{\Var}{\operatorname{Var}}
\newcommand{\esssup}{\operatorname*{ess\,sup}}
\newcommand{\essinf}{\operatorname*{ess\,inf}}

\title{On the propagation factor $F_\star$: why it is not the quantity to report}
\author{}
\date{2026-08-02}

\begin{document}
\maketitle

\begin{abstract}
This report examines the propagation factor $F_\star$ introduced in the red
passages of \texttt{F\_star/main\_F\_star.tex} and asks three questions: whether
it can serve as the error bound the experiments report, what its evaluation
costs relative to the reported factor $\hat F_\Sigma$, and whether it resolves
the difficulty that a singular molecular potential creates for the supremum
factor $F_{\Sigma+}$. The theoretical content of $F_\star$ is sound and its
central observation is correct: the exponential dependence on the stage count
$K$ in the published constant is an artifact of the recursion, and telescoping
over the injected errors removes it. What follows is not a criticism of that
argument. The conclusion is that $F_\star$ is a good remark and a poor
statistic. It is governed by a two-sided norm, and the side it adds --- an
essential infimum --- is precisely the side a singular potential destroys, so
it does not settle the singularity difficulty but inherits it twice. Its
sample counterpart is an exponentiated extremum, which has no limit in the
particle count, is a property of the realized genealogy rather than of the
stage schedule, and is biased low at exactly the stages that dominate it.
Finally, it cannot be computed from what the completed runs saved. Section~8
records the decision taken on the strength of this analysis: $F_\star$ is not
mentioned anywhere in the paper.
\end{abstract}

\section{What $F_\star$ is}

Write $\mu_t\propto\exp(-U_t)$ with $U_t=(1-t)U_0+tU$ for the interpolation
from the source to the target, $\mu_k:=\mu_{t_k}$ for the accepted stage
targets, $\nu_k$ for the stage proposal, and $w_k:=\D\mu_k/\D\nu_k$ for the
stage weight. After the change of variables that removes the stage maps, the
accumulated ratio
\begin{equation}
    R_k:=\prod_{j=k+1}^{K}r_j=\frac{\D\varrho_K}{\D\varrho_k},
    \qquad R_K:=1,
    \label{eq: accumulated}
\end{equation}
carries the stage-$k$ target all the way to the endpoint in one reweighting.
The proposed theorem assumes that each $R_k$ is bounded on both sides off
null sets of the source,
\begin{equation}
    0<\underline R_k:=\essinf R_k \Le \esssup R_k=:\overline R_k<\infty,
    \qquad
    \varrho^\star_k:=\frac{\overline R_k}{\underline R_k}\Ge1,
    \label{eq: two-sided}
\end{equation}
and concludes the same bound as the published theorem with the factor
\begin{equation}
    F_\star:=\sum_{k=0}^{K}\varrho^\star_k
    \qquad\text{in place of}\qquad
    F_{\Sigma+}=\sum_{k=0}^{K}\ \prod_{j=k+1}^{K}3\lVert w_j\rVert_\infty .
    \label{eq: two-factors}
\end{equation}

Two things are genuinely gained, and they should be stated before the
objections. First, the geometric factor $3^{K-k}$ disappears. In the neutral
case $\mu_k\equiv\mu$ with identity stage maps, every $R_k\equiv1$ and
$F_\star=K+1$, whereas $F_{\Sigma+}=(3^{K+1}-1)/2$ and the constant grows like
$9^{K}$ although nothing whatever is lost by a stage. Second, the supremum of a
product is replaced by the spread of a single accumulated ratio, and successive
stage weights need not peak at the same configurations, so the replacement is
typically far from an equality. The accompanying ordering lemma also shows that
$F_\star$ dominates the accumulated factor $F_{\mathrm{acc}}$ and bounds the
asymptotic variance \emph{linearly} rather than through its square. All of this
is correct and worth recording.

The rest of this report concerns whether the resulting quantity can be
\emph{reported}.

\section{The hypothesis is two-sided, and the new side is the fragile one}

The published theorem assumes only $\lVert w_k\rVert_\infty<\infty$. The
proposed theorem assumes \eqref{eq: two-sided}, which is strictly stronger in
kind, not merely in degree: it adds a lower bound on a density. The red text
acknowledges this --- ``the bound is vacuous when a stage weight comes
arbitrarily close to zero on the carried source support'' --- and states
plainly that ``the two results are therefore not ordered in general''.

That non-ordering is decisive for a reported quantity. A factor whose
finiteness is not implied by the finiteness of the factor it replaces cannot
be substituted for it in the experimental tables: there are schedules on which
$F_{\Sigma+}<\infty$ and $F_\star=\infty$. A reported column must be defined
for every run in it.

It is worth being explicit about where the lower bound fails, because for the
molecular targets it fails for a structural reason rather than an accidental
one. Taking identity stage maps for clarity, the accumulated ratio of
\eqref{eq: accumulated} is
\begin{equation}
    R_k \ \propto\ \frac{\mu_K}{\mu_k}
    \ =\ \exp\bigl(-(1-t_k)(U-U_0)\bigr),
    \label{eq: R-explicit}
\end{equation}
so the two ends of $R_k$ are governed by the two ends of the energy difference
$U-U_0$:
\begin{itemize}
    \item $\underline R_k=0$ wherever $U-U_0\to+\infty$, that is, in the
    repulsive core where the target potential diverges;
    \item $\overline R_k=\infty$ wherever $U-U_0\to-\infty$, that is, in the
    far field where the Gaussian source potential diverges while the target
    potential does not.
\end{itemize}
The first item is exactly the singular-potential difficulty. The second is
present for every stage $k<K$ whenever the source is Gaussian on an unbounded
component and the target energy is capped. Both ends therefore fail for the
same schedules on which $\lVert w_k\rVert_\infty$ was already in doubt, and the
lower end fails for a reason the upper end did not have.

\section{Does $F_\star$ settle the singularity error?}

It does not. The relevant comparison is between norms. The manuscript's
Appendix B.3 already explains why $F_\Sigma$ survives a singular potential and
$F_{\Sigma+}$ does not: an effective sample size is an $L^2$ quantity and a
supremum is an $L^\infty$ quantity, and a square-integrable weight can have
$\ESS(w_k)>0$ while $\lVert w_k\rVert_\infty=\infty$. The ordering of the three
candidates by fragility is therefore
\begin{equation*}
    \underbrace{L^2}_{\textstyle F_\Sigma}
    \ \prec\
    \underbrace{L^\infty}_{\textstyle F_{\Sigma+}}
    \ \prec\
    \underbrace{L^\infty\ \text{and}\ L^{-\infty}}_{\textstyle F_\star},
\end{equation*}
and $F_\star$ sits at the fragile end. It asks for more control of the weight
than the factor that already fails, so it cannot repair a failure caused by
insufficient control.

The point survives even when both ends are finite. Suppose the domain is
compact --- a torus component, say --- or that the regularized potential
$U^{\{\rho\}}$ caps the energy at the threshold $e$, so that
\eqref{eq: two-sided} holds. Then \eqref{eq: R-explicit} gives
\begin{equation}
    \varrho^\star_k
    \ =\ \exp\Bigl((1-t_k)\bigl[\,\esssup(U-U_0)-\essinf(U-U_0)\,\bigr]\Bigr)
    \ =\ \exp\bigl((1-t_k)\,\osc(U-U_0)\bigr).
    \label{eq: rho-star-energy}
\end{equation}
The summand is the exponential of an accessible energy range measured in units
of $k_{\mathrm B}T$. For the molecular targets that range is not a small
number: the regularization threshold $e$ alone is a large multiple of
$k_{\mathrm B}T$, and the source and target potentials disagree by that order
over the sampled region. A factor of the form $\exp(\text{hundreds})$ is not a
diagnostic that a reader can compare between two runs; it is a number whose
logarithm is the only readable part.

So the regularization, not $F_\star$, is what makes the constant finite, and
what it leaves is finite in the same sense that $\exp(300)$ is finite. This is
the precise sense in which $F_\star$ does not settle the singularity error: it
converts an infinite factor into an astronomically large one, by borrowing the
cap that the sampling potential already applies, and it does so only on the end
where the cap acts.

\section{The sample counterpart is an exponentiated extremum}

Even granting the hypothesis, the reported quantity would not be $F_\star$ but
its sample counterpart. With $S_k^i$ the accumulated suffix log-weight along
the lineage of particle $i$, the proposal is
\begin{equation}
    \hat F_\star=\sum_{k=0}^{K}\exp\Bigl(\max_i S_k^i-\min_i S_k^i\Bigr).
    \label{eq: F-star-hat}
\end{equation}
Four properties of \eqref{eq: F-star-hat} disqualify it as a reported number.

\paragraph{It has no limit in $N$.} A maximum and a minimum over $N$ draws are
extremum statistics. For a log-weight with unbounded support the observed range
$\max_iS_k^i-\min_iS_k^i$ grows without bound as $N$ increases, so
$\hat F_\star$ increases with the particle count and does not converge. It
therefore measures the sample size as much as it measures the run. By contrast
$\ESS(\mathbf W_j)$ is a second moment: it converges as $N$ grows, and
$\hat F_\Sigma$ inherits that stability. The manuscript reports factors
computed at $N$ between $2.4\times10^5$ and $3.2\times10^5$ and compares them
across molecules; a statistic that drifts with $N$ cannot be used that way.

\paragraph{The drift is then exponentiated.} The quantity that grows with $N$
sits inside an exponential. An increase of the observed log-range by $5$, which
a tenfold increase in $N$ can easily produce for a heavy-tailed weight,
multiplies the reported factor by about $150$. The reported value is thus not
merely unstable but unstable on a multiplicative scale that swamps any
difference between two training objectives, which is the comparison the
factor exists to make.

\paragraph{It is a property of the genealogy, not of the schedule.} The
accumulated ratio $\hat R_k^i=\prod_{j>k}W_j^{a_j(i)}$ is taken along the
ancestral line $a_j(i)$ produced by the realized resampling draws. Two
inference passes with the same trained flows, the same accepted schedule and
different resampling seeds produce different lineages and therefore different
$\hat F_\star$. The published factors do not have this property:
$\hat F_\Sigma$ and $\hat F_{\Sigma+}$ are functions of the stage weight
vectors alone. A factor that characterizes ``one accepted schedule'', which is
how the manuscript uses the term, must not depend on the resampling draws of a
particular replay.

\paragraph{Its bias runs the wrong way.} The red text records the defect
itself: resampling makes lineages coalesce, so the number of distinct
ancestors at stage $k$ falls as $K-k$ grows, and the observed spread of
$\hat R_k^i$ for early $k$ is computed over few distinct values and is biased
low. The early stages are precisely the ones with the largest accumulated
ratio and therefore the largest summands in \eqref{eq: F-star-hat}. The
estimator is thus optimistic exactly where it matters most, and its bias
cannot be signed away by increasing $N$, since coalescence is driven by the
number of stages.

\section{The calculation is more expensive, and the saved runs cannot support it}

The cost argument is concrete rather than asymptotic.

\paragraph{What the reported factors need.} $\hat F_\Sigma$ and
$\hat F_{\Sigma+}$ are read off the stage weight vectors
$\mathbf W_j=(W_j^1,\dots,W_j^N)$, which every run evaluates anyway: one
second moment and one maximum per vector, a single pass over data already on
disk. The completed chiral runs persist exactly these, per stage:
\texttt{flow\_log\_weights.npy}, \texttt{sharpen\_log\_weights.npy},
\texttt{samples.npy}, the two acceptance arrays, and \texttt{metadata.json}.

\paragraph{What $F_\star$ needs, and what is missing.} Equation
\eqref{eq: F-star-hat} needs the resampling ancestor indices $a_j(i)$. No
ancestor, index or lineage array appears anywhere under the chiral run
directories. The saved artifacts therefore do not determine $\hat F_\star$,
and no amount of post-processing recovers it: the genealogy was not recorded
and cannot be reconstructed from weights and positions. Reporting $F_\star$
for the completed runs would require re-running inference on all three
molecules with ancestry recorded.

\paragraph{The coalescence-free alternative is a second training-scale replay.}
The remedy offered for the coalescence bias is to push each stage-$k$
validation particle through the remaining maps $G_{k+1}^{-1},\dots,G_K^{-1}$
without resampling. This is $O(K)$ flow evaluations per particle per starting
stage and $O(K^2N)$ in total. For alanine dipeptide at $K=10$ and
$N=3.0\times10^{5}$ that is on the order of $10^{7}$ flow evaluations, against
the $O(KN)\approx3\times10^{6}$ of a single inference replay: a factor of $K$
more work than the pass that produced the samples in the first place, to
obtain one column of one table.

\paragraph{And the summary side is also unavailable.} The per-stage summaries
saved with the reported factors record the effective sample size and the
largest rescaled weight, but not the smallest. The infimum side of
\eqref{eq: two-sided}, which is the side $F_\star$ adds, was never persisted
for any of the three molecules.

\section{What the measured runs show}

Table~\ref{tab: measured} reproduces the factors computed from the three
completed chiral runs. These are measured values from the saved artifacts, not
estimates.

\begin{table}[htb]
\centering
\begin{tabular}{lrrrrr}
\toprule
molecule ($d$) & $N$ & stages & $\hat F_\Sigma$ & $\hat F_{\Sigma+}$ & $\log_{10}\hat F_{\Sigma+}$\\
\midrule
$(2R,3R)$-2,3-butanediol (42) & $2.4\times10^5$ & 8  & 37.45  & $1.01\times10^{20}$ & 20.01\\
\textsc{l}-alanine dipeptide (60) & $3.0\times10^5$ & 10 & 68.08  & $3.55\times10^{26}$ & 26.55\\
Ac-Pro-NHMe (72)              & $3.2\times10^5$ & 11 & 176.75 & $3.37\times10^{31}$ & 31.53\\
\bottomrule
\end{tabular}
\caption{Propagation factors of the three completed chiral runs, accumulated
over the selected flow or identity weight of every accepted stage followed by
its sharpening weight. Reproduced from the saved run summaries.}
\label{tab: measured}
\end{table}

Two features of these runs bear on the question. First, the spread between the
two published factors is already $18$ to $29$ orders of magnitude, and the
whole of that spread is the product structure that $F_\star$ is designed to
remove. The diagnosis in the red text is therefore correct on the real runs and
not only on the six-state illustration: the exponential growth is an artifact.
Second, the per-stage maxima that generate it are individually modest --- the
largest rescaled weight over all entries of all three runs is $106.07$, at
alanine dipeptide stage 2 --- while the corresponding effective sample sizes
never fall below $0.410$. A schedule whose worst stage retains $41\%$ of its
particles is not a pathological one, and yet its supremum factor is $10^{26}$.
This is the strongest available evidence that the exponential constant should
be read as an artifact, and it is an argument for the remark, not for a new
reported column.

\section{Conclusion and recommendation}

\begin{enumerate}
    \item \textbf{The mathematics of $F_\star$ is correct and its diagnosis is
    right.} The telescoping identity, the theorem, and the ordering lemma all
    hold as stated, and the exponential dependence on $K$ in the published
    constant is genuinely an artifact of re-entering a worst case at every
    stage.
    \item \textbf{It is not a reportable error bound.} Its hypothesis is
    two-sided and is not implied by the hypothesis it replaces, so the two
    theorems are not ordered and $F_\star$ can be infinite where $F_{\Sigma+}$
    is finite.
    \item \textbf{It does not settle the singularity difficulty.} It requires
    control of the weight in $L^\infty$ \emph{and} of its reciprocal, so it is
    strictly more fragile than the factor that already fails on a singular
    potential. Where the regularized potential does make it finite, the value
    is $\exp(\osc(U-U_0))$ by \eqref{eq: rho-star-energy}, an exponential of an
    energy range in units of $k_{\mathrm B}T$.
    \item \textbf{Its estimator is unsuitable independently of all that.} It is
    an exponentiated extremum: it drifts upward with $N$, it depends on the
    realized resampling genealogy rather than on the accepted schedule, and its
    coalescence bias is low at the dominant stages.
    \item \textbf{It costs a second replay, and cannot be evaluated at all on
    the completed runs}, which persisted neither ancestor indices nor the
    minimum-weight side of the per-stage summaries.
\end{enumerate}

The reported factor should therefore remain $\hat F_\Sigma$, which is bounded
below by $K+1$, built from second moments that converge in $N$, computable from
vectors already on disk, and independent of the resampling draws. The
manuscript's existing position --- that $\hat F_\Sigma$ is not an error bound but
is the honest summary of accumulated stagewise weight degeneracy along one
accepted schedule --- is the one the evidence supports.

\section{Decision}

\noindent\textbf{Decided by Xuda Ye, 2026-08-02: $F_\star$ is not to be
mentioned at all, anywhere in this paper.}

The decision is one of scope rather than of correctness. This is a
computational frontier paper: it proposes a training objective and a generator,
and it reports what a run costs and what it recovers. The propagation factor
enters it as a summary of accumulated stagewise weight degeneracy, not as the
subject of the paper. A second factor, a second theorem, its two-sided
hypothesis, the ordering lemma relating the two, and the discussion of which is
honest in which direction would add a strand of analysis that the paper's
readership is not reading it for, and that no experiment in it uses. The
material in \texttt{main\_F\_star.tex} is therefore held outside the
manuscript in its entirety, including the parts of it that are correct.

Three consequences are worth recording so that the decision is not partially
undone later.

\begin{itemize}
    \item \textbf{Nothing in the manuscript needs to change.} $F_\star$,
    $\varrho^\star_k$, $F_{\mathrm{acc}}$, the accumulated ratio and the
    two-sided hypothesis appear nowhere in \texttt{Paper/main.tex}. The
    symbols $R_k$ and $\essinf$ that do occur there are a different object and
    a different use: $R_k^i=r_k(Z_{k-1}^i)$ is the reduced-procedure stage
    weight of Appendix B, and $\essinf$ appears only in the definition
    $\osc(\varphi)=\esssup\varphi-\essinf\varphi$ in Section 4.4.
    \item \textbf{The pessimism in $K$ is already stated without $F_\star$.}
    Appendix B.3 makes that point on its own terms, with the neutral-case
    computation $F_{\Sigma+}=(3^{K+1}-1)/2$ and a constant growing like
    $9^{K}$, and concludes that the theorem establishes the rate $N^{-1/2}$ and
    the finiteness of the constant rather than its size. That paragraph needs
    no reference to a second factor, and none should be added to it.
    \item \textbf{No column, remark, footnote or appendix.} The exclusion
    covers the reported tables, the discussion of what the factor does and does
    not measure, and any forward-looking sentence promising a sharper constant.
    A partial mention would carry the cost the decision is avoiding --- a
    reader asking what the other factor is --- without the benefit of the full
    treatment.
\end{itemize}

The analysis in Sections 1--7 above is retained here, and only here, as the
record of why the question was asked and how it was answered. Should a later
and more theoretical paper take up the constant itself, this report and
\texttt{main\_F\_star.tex} are the starting point: the telescoping argument,
the theorem, and the ordering lemma are sound, and the objections in
Sections 2--6 are objections to reporting $F_\star$ in \emph{this} paper, not
to the result.

\end{document}
